% Full midpoint concatenation and finite-family proof of Kamp III.2, Lemma 5.

\section{Kamp's composition through an interior point}
\label{sec:annex-p09-kamp-iii2-lemma5}

\begin{annexitem}{annex:p9.kamp.elementary.midpoint}{thm:p09-kamp-elementary-midpoint}{Concatenation with all three midpoint labels retained}
The source reduces display~(21) to two basic patterns and combines their
shared endpoint labels by proper substitution
\cite[ch.~III, p.~79]{Kamp1968}.

\begin{proof}
Fix a nonempty strict linear order with arbitrary unary interpretations,
giving a structure \(\mathcal M\). Use the labelled realization
equivalence from \cref{thm:p09-kamp-elementary-intervals}.
For a label with free parameter \(w\), saying that it holds at \(t\)
means satisfaction in \(\mathcal M\) with \(w\) assigned that point.
For an interval formula, its two displayed point arguments mean
assignment to its free endpoints. All remaining variables are bound.

\paragraph{The joined word and its elementary presentation.}
The words \(W,V\) have \(a,b\geq1\) intervals. The joined word
retains their \(a+b\) interval positions and identifies the last left
point position with the first right point position. Its number of
points is therefore \((a+1)+(b+1)-1=a+b+1\), and its total length
is \(2(a+b)+1\). At the identified position the label is
\((\eta\land\alpha_a)\land\gamma_0\). All three operands belong to
\(\operatorname{E}(w)\); two applications of that class's conjunction
constructor place the new label in \(\operatorname{E}(w)\).
Every other label was already in that class.

Thus this is an elementary labelled word with \(a+b\geq2\)
intervals. The preceding elementary-interval presentation theorem gives
an actual \(\operatorname{Ep}(x,y)\) presentation. More explicitly,
choose a basic pattern in \(\operatorname{Bp}_{a+b}(x,y)\) with
distinct placeholder letters for its \(2(a+b)+1\) occurrences,
all absent from every label, then substitute each label properly at
its own occurrence. The source-list construction proved in
\cref{annex:p9.kamp.elementary.intervals} supplies the finite
initial-segment substitution list if needed. It uses no closure under
semantic equivalence as a syntactic membership rule.

\paragraph{An existential midpoint yields a joined realization.}
Let the outer endpoints be \(s<t\), and suppose the existential
formula in the lemma holds there. It supplies a point \(d\) with
\(s<d<t\), at which \(\theta\) holds, and at which the two interval
formulas hold on their respective endpoint pairs \((s,d)\) and
\((d,t)\). Their elementary words therefore have realizations
\[
      s=u_0<\cdots<u_a=d,\qquad
      d=v_0<\cdots<v_b=t.
\]
The replacement of the free parameter in \(\theta\) gives
\(\eta(d)\) exactly when \(\theta\) holds with \(z=d\), by the
free-variable substitution identity proved in
\cref{annex:p9.kamp.proper.substitution}.

Join the sequences into
\[
               u_0,\ldots,u_{a-1},d,v_1,\ldots,v_b.
\]
This sequence is strictly increasing. The only shared point is the
displayed \(d=u_a=v_0\): every earlier left point lies below \(d\)
and every later right point above it. Its point labels through position
\(a-1\) are \(\alpha_0,\ldots,\alpha_{a-1}\), supplied by the
left realization. At position \(a\), the left realization supplies
\(\alpha_a(d)\), the right supplies \(\gamma_0(d)\), and the
midpoint hypothesis supplies \(\eta(d)\). These are exactly the
three conjuncts of its new label. At position \(a+j\), for
\(1\leq j\leq b\), the right realization supplies \(\gamma_j(v_j)\).

For interval index \(i\), \(1\leq i\leq a\), the joined interval
is the original \((u_{i-1},u_i)\) and hence satisfies \(\beta_i\)
throughout. For index \(a+j\), \(1\leq j\leq b\), it is the
original \((v_{j-1},v_j)\) and satisfies \(\delta_j\) throughout.
These exhaust the intervals; there is no interval crossing past the
shared point without stopping there. Thus the joined sequence realizes
the entire word \(W\mathbin{\star_\eta}V\) at endpoints \(s,t\).
The realization equivalence proves the required joined pattern.

\paragraph{A joined realization supplies the same midpoint to both pieces.}
Conversely, suppose the joined pattern holds at \(s<t\). Choose a
realization
\[
                         s=r_0<\cdots<r_{a+b}=t.
\]
Set \(d=r_a\). Since \(a\geq1\) and \(b\geq1\),
\(0<a<a+b\), so \(s<d<t\). Its point label gives separately
\(\eta(d)\), \(\alpha_a(d)\) and \(\gamma_0(d)\).

The prefix \(r_0<\cdots<r_a\) realizes \(W\). Its point labels
before the last point are unchanged; the last required label
\(\alpha_a\) is a conjunct of the combined label. Its intervals
at indices \(1,\ldots,a\) have their original \(\beta_i\) labels.
The realization equivalence therefore yields \(\rho\) at \((s,d)\).
The suffix \(r_a<\cdots<r_{a+b}\) realizes \(V\): its first
required label \(\gamma_0\) is another conjunct of the combined
label, its later point labels are the original \(\gamma_j\), and
its \(b\) intervals have their original \(\delta_j\) labels.
Thus \(\sigma\) holds at \((d,t)\). Parameter substitution transfers
\(\eta(d)\) back to \(\theta\) with \(z=d\).
The same point \(d\) consequently satisfies every conjunct of the
guarded existential; different witnesses for the two sides have not
been substituted for it.

These arguments prove the displayed equivalence on every ordered
endpoint pair in every linear structure. If the outer points are not
ordered, its guard is false. In fact both formulas themselves are
false there: the existential requires \(x<z<y\), and the joined
basic skeleton has at least two intervals, so its outer recursive
constructor already requires a strict interior witness. The stated
guarded identity needs only the ordered case.

\paragraph{The exact source endpoint-label substitution.}
Here is the connection with the source's formula \(\pi'\).
First take pure basic skeletons for \(\rho,\sigma\), with disjoint
placeholder letters at all their occurrences. The elementary-interval
history theorem supplies these skeletons even if the original formulas
shared letters or used the same label at several positions. Choose all
placeholders absent from all eventual labels, and choose another fresh
unary letter, written \(Q_0\) as in the source, different from \(P_0\)
and absent from both skeletons and all eventual labels. This \(Q_0\) denotes a chosen
predicate letter of \(L\); it is not the fixed letter \(P_0\)
used in the source tautology convention.

Let \(Q'\) be the last predicate occurrence in the left skeleton's
word, and let \(Q''\) be the first in the right skeleton's word.
They are different letters by the disjoint choice. Remove that first
\(Q''\) from the right word, and append the remaining right word
to the complete left word. The result is a basic word with \(a+b\)
intervals. Choose \(\pi\in\operatorname{Bp}_{a+b}(x,y)\) having
that word, and form the proper substitution
\[
 \pi'\defeq
 \operatorname{psub}_{Q'}
   \bigl(\pi,\,(Q_0(z)\land Q'(z))\land Q''(z)\bigr).
\]
Here \(\operatorname{psub}_{Q'}\) means the previously defined
operation at the index of the unary letter \(Q'\).
The replacement is an \(\operatorname{At}(z)\) formula, hence belongs
to \(\operatorname{E}(z)\). Therefore \(\pi'\) is in
\(\operatorname{Ep}(x,y)\). The unique original occurrence of
\(Q'\) is the midpoint position of \(\pi\). Its substitution
places \(Q_0,Q',Q''\) together at that point. The newly inserted
\(Q'\) is retained by the single-pass substitution convention;
it is not expanded repeatedly. No other position of \(\pi\) is
changed by this step.

If the bound midpoint variable of \(\pi\) is called \(v\), the
inserted formula at that position is
\((Q_0(v)\land Q'(v))\land Q''(v)\). The replacement's free
parameter \(z\) is replaced by the argument \(v\), and the existing
quantifier of \(\pi\) binds those argument occurrences. This is
predicate substitution at an atom, not substitution into a bound
variable declaration. The symbol \(z\) does not become a new free
endpoint. Other bound variables are fresh, as required by proper
substitution.

The two realization arguments above, with the elementary labels
specialized to these distinct letters, prove that \(\pi'\) agrees
with source display~(21): a strict midpoint satisfying \(Q_0\),
with the left skeleton ending there and the right skeleton beginning
there. In particular deleting \(Q''\) from the concatenated word
does not delete its condition: the proper substitution restores it
in the midpoint conjunction.

Now restore the elementary label attached to each skeleton placeholder,
including \(Q'\) and \(Q''\). The fresh-letter choice ensures
that later such substitutions do not change a previously inserted
label. The proved finite-history construction treats any substitutions
already present in the original \(\rho,\sigma\) in their actual
order. Then substitute \(\theta(z)\) for \(Q_0\).
The other labels contain no \(Q_0\), so they are unchanged.
The midpoint label becomes precisely
\((\eta\land\alpha_a)\land\gamma_0\) at its current argument.
All other point and interval labels are those of the joined word.
Each step preserves the elementary interval class by R6, proved in
\cref{annex:p9.kamp.elementary.intervals}.

For the corresponding property-\(I\) reduction used in the source,
R7 in \cref{thm:p09-kamp-interval-substitution} permits the same
proper substitutions in a guarded interval equivalence on
\(\mathsf{Com}\). The substitution for a fresh \(Q_0\) in
\(\exists z\,(x<z<y\land Q_0(z)\land\rho\land\sigma)\) changes
exactly its midpoint atom to \(\theta\); its two strict inequalities
and its other constituents are unchanged. This verifies explicitly
the reduction from a general elementary midpoint condition to a
fresh unary atom and back. The argument uses the established
substitution semantics and class closure, with their declared domains.

Every operation is finite. The complete word retains every input
interval label and every endpoint label, using conjunction at the
identified point. Its labelled formula has exactly \(x,y\) free,
by the elementary presentation theorem; proper substitutions preserve
that set. Neither a density premise nor completeness of the order
entered either realization direction.
\end{proof}
\end{annexitem}

\begin{annexitem}{annex:p9.kamp.iii2.lemma5}{thm:p09-kamp-iii2-lemma5}{All finite side conjunctions, disjunct combinations and midpoint cases}
The source's Lemma~5 occupies printed pp.~78--79. We use the given
interval representations on \(\mathsf{Com}\), exactly as in its
displayed premises on p.~78
\cite[ch.~III, pp.~78--79]{Kamp1968}.

\begin{proof}
Fix the three finite families in the body, with the indicated free
variables, and rename all their bound variables fresh before placing
them under the one outer midpoint quantifier. First construct the
three aggregate conditions and their side representations.
In semantic assertions, point arguments such as \(\Gamma(s,t)\),
\(\Phi(s,d)\), \(\Psi(d,t)\) and \(\theta(d)\) mean truth with
the displayed free endpoints assigned those points.

\paragraph{The elementary midpoint conjunction.}
For \(m=1\), \(\theta\) is exactly \(\chi_1\).
For \(m>1\), initialize a conjunction with \(\chi_1\), then conjoin
each \(\chi_i\) in increasing index order. At every step both
operands lie in \(\operatorname{E}(z)\), so its binary conjunction
closure gives a member of that same class. Induction through the
finite list proves \(\theta\in\operatorname{E}(z)\).
For \(m=0\), the explicitly chosen \(T^K_z\) already belongs to
\(\operatorname{At}(z)\subseteq\operatorname{E}(z)\).
The parameter-renaming result in
\cref{annex:p9.kamp.elementary.intervals} therefore places
\(\eta=\theta[w/z]\) in \(\operatorname{E}(w)\) in every case.

\paragraph{The left side and every choice of its elementary disjuncts.}
For \(n=1\), use the supplied finite elementary representation of
\(\phi_1\). Suppose a current list represents
\(\phi_1\land\cdots\land\phi_{j-1}\) under the guard \(x<z\)
on complete orders. The given list for \(\phi_j\) represents it
under the same guard and on the same class. Distributing these two
finite disjunctions gives every pair of their disjuncts. The preceding
Lemma~4 replaces each pair by its full finite fusion list in
\(\operatorname{Ep}(x,z)\). Concatenate all those lists. On
\(\mathsf{Com}\), under \(x<z\), the new disjunction is therefore
equivalent to \(\phi_1\land\cdots\land\phi_j\).

This proves the induction step and yields the stated list
\(\mathcal L\) after the last input. It also proves completeness
of the enumeration: at the first stage every first-input disjunct is
listed. At the next stage every current choice is paired with every
choice for the next input, and every fusion is retained. Induction
therefore includes every original tuple of disjunct choices and every
fusion branch produced along that tuple. No one disjunct or one
fusion is chosen to represent the entire conjunction.

For \(n=0\), \(\Phi(x,z)=T^K_x\) is true for every assignment.
The proposed neutral elementary pattern
\[
                 \mathcal B_p((T^K_w,T^K_w,T^K_w);x,z)
\]
is also true whenever \(x<z\): its two point labels are tautologies,
and every point of its open gap satisfies its tautological interval
label. It has a one-interval elementary presentation, not an added
zero-interval basic stage. Its one-member list consequently represents
\(\Phi\) on ordered endpoints, including when those endpoints are
adjacent. We obtain in all cases
\[
 \mathsf{Com}\models x<z\longrightarrow
              \left(\Phi(x,z)\longleftrightarrow
                                      \bigvee_{\rho\in\mathcal L}\rho\right).
\]

\paragraph{The right side and all its disjunct choices.}
For \(p=1\), keep the supplied representation of \(\psi_1\).
For \(p>1\), at each stage distribute the current representation
against every disjunct of the next \(\psi_k\), now at endpoints
\(z,y\). Lemma~4 supplies every elementary fusion in
\(\operatorname{Ep}(z,y)\). The induction assertion at stage \(k\)
is equivalence, on \(\mathsf{Com}\) under \(z<y\), to
\(\psi_1\land\cdots\land\psi_k\). Its base is the supplied first
representation and its step is precisely the guarded distribution and
fusion identity, so induction proves the assertion for the final list
\(\mathcal R\). The same enumeration induction includes every original
right-side choice and every intermediate fusion branch.

For \(p=0\), \(\Psi(z,y)=T^K_y\) is true everywhere. The neutral
pattern
\(\mathcal B_p((T^K_w,T^K_w,T^K_w);z,y)\)
is true on every \(z<y\), by its two point tautologies and its
universal interval tautology. Thus its one-member list represents
\(\Psi\). In every case,
\[
 \mathsf{Com}\models z<y\longrightarrow
              \left(\Psi(z,y)\longleftrightarrow
                                    \bigvee_{\sigma\in\mathcal R}\sigma\right).
\]

If any given representation is written as an empty disjunction,
replace it first by the false elementary pattern allowed by
\cref{def:p09-kamp-interval-property}. This is a valid replacement
under the same ordered-endpoint guard and yields a nonempty list.
The preceding inductions apply to those lists. Their disjunctions
may still be false; nonemptiness of a syntactic list does not assert
existence of a realization.

\paragraph{Distribution under the one guarded existential.}
Fix a structure \(\mathcal M\) over a complete linear order and
outer points \(s<t\). Suppose \(\Gamma(s,t)\) holds. Choose its
witness \(d\) with \(s<d<t\). This one assignment gives both
\(\Phi(s,d)\) and \(\Psi(d,t)\), as well as \(\theta(d)\).
The left guarded equivalence gives at least one
\(\rho\in\mathcal L\) true at \((s,d)\), since \(s<d\).
The right guarded equivalence gives at least one
\(\sigma\in\mathcal R\) true at \((d,t)\), since \(d<t\).
The same \(d\) witnesses
\(\exists z\,(x<z<y\land\theta\land\rho\land\sigma)\) for this pair.
The choices of disjuncts may depend on \(d\); the finite disjunction
includes all of them.

Conversely, if some pair \(\rho,\sigma\) has such a midpoint
witness \(d\), its strict inequalities permit both side equivalences
to be applied in their reverse directions. They give
\(\Phi(s,d)\) and \(\Psi(d,t)\). Together with \(\theta(d)\),
these prove \(\Gamma(s,t)\). We have therefore established
\[
\begin{split}
 \Gamma(s,t)\quad\Longleftrightarrow\quad
 \bigvee_{\rho\in\mathcal L}\ \bigvee_{\sigma\in\mathcal R}
   \exists d\,\bigl(&s<d<t\land\theta(d)\\
                 &{}\land\rho(s,d)\land\sigma(d,t)\bigr).
\end{split}
\]
Here the displayed point arguments denote satisfaction under the
corresponding assignments, not additional parameters of an input formula.
This proof keeps one witness for all midpoint conjuncts. It does not
replace the conjunction by separate existential witnesses for the two
sides. The side equivalences need only hold on ordered pairs because
the outer witness guard supplies both required orders.

\paragraph{Every pair becomes one concatenated elementary pattern.}
For a selected pair, obtain the labelled words
\(W_\rho,V_\sigma\) and use the preceding annex item with midpoint
label \(\eta\). It gives an
\(\operatorname{Ep}(x,y)\) presentation for
\[
             \mathcal B_p(W_\rho\mathbin{\star_\eta}V_\sigma;x,y)
\]
and its equivalence to that pair's midpoint existential. The identity
was proved on every linear order from the actual elementary
realizations, so it holds on the complete order \(\mathcal M\).
Replace every pair by its resulting formula in the displayed finite
disjunction. The previous paragraph proves exactly the body's
\(\mathsf{Com}\) equivalence for \(\Gamma\) under \(x<y\).
This is property \(I(x,y)\).

The given arbitrary \(I\)-formulas have been replaced only on complete
orders. No step applies their representation equivalences to an
arbitrary incomplete order. The elementary concatenation identity's
larger order scope came from its separate realization proof, not from
changing the scope of these premises.

The free-variable sets are preserved as required. Each original
midpoint formula has exactly \(z\) free; each left and right formula
has only its stated endpoints possibly free. After bound renaming,
the outer existential binds their free \(z\)-occurrences, leaving
no variable other than \(x,y\). Both \(x,y\) occur freely in the
strict guard even if a family is empty or a constituent is a sentence.
Each constructed disjunct has exactly \(x,y\) free by the elementary
presentation theorem. Its internal midpoint is bound at the joined
word's specified interior position.

\paragraph{Single families, false conditions and vacuous intervals.}
If a family contains just one formula, its conjunction is that formula
and the corresponding finite-list construction performs no unnecessary
conjunction step. Repeated input formulas or repeated disjuncts cause
no difficulty: the same truth may be required more than once, and
duplicate alternatives may remain in the list.

If some \(\chi_i\) is false at every point, then \(\theta\) is
false at every possible midpoint. Each joined word has \(\eta\)
as a conjunct of its central point label, so no output disjunct has
a realization. Both sides of the claimed equivalence are false.
If some \(\phi_j\) is false on every ordered pair, the aggregate
left side is false under \(x<z\), and its proved representation
has no realizing disjunct there. Every joined realization would
restrict to such a left realization by the preceding annex item,
so no output disjunct is realizable. The same argument, using the
right suffix, handles a \(\psi_k\) false on every ordered pair.
These observations include a supplied empty disjunction replaced
by the false elementary pattern.

Falsehood of an interval \emph{label} alone has a different effect.
The universal condition on an empty open gap is true even if its
label is false everywhere. Thus such a label can require that a
gap contain no point. For example, on \(0<1<2\), use the word
\((T^K_w,F^K_w,T^K_w)\) on each side and \(T^K_w\) at the
midpoint. The left realization \((0,1)\) and right realization
\((1,2)\) satisfy their false interval labels vacuously. Their
joined realization \((0,1,2)\) satisfies both false gap labels.
The composition is true at \((0,2)\). No proof step infers the
falsehood of a whole pattern merely from a false gap label.

Contradictory requirements \emph{at the shared point} do make the
composition false. For example, if the last left label is \(P(w)\)
and the first right label is \(\neg P(w)\), the combined point
label contains \(P(w)\land\neg P(w)\), regardless of the midpoint
condition. It has no satisfying point. Both labels must hold at
the same witness, exactly as the source's conjunction
\(Q_0(z)\land Q'(z)\land Q''(z)\) requires.

\paragraph{Empty families and the strict witness condition.}
For each zero family size, the assigned tautology is a formula of the
declared source grammar and is true in every structure. Therefore
replacing an absent family by that formula adds no semantic requirement.
The neutral side patterns have one interval whose endpoints are the
already required outer point and midpoint; they do not require another
point inside that interval. This verifies separately the cases
\(m=0\), \(n=0\) and \(p=0\), as well as every combination of them.

When all three sizes are zero, \(\Gamma(x,y)\) is equivalent to
\(\exists z\,(x<z<y)\), because all three aggregate formulas are
tautologies. Both neutral side words have one interval, so the output
is the word
\[
        (T^K_w,T^K_w,\,
          (T^K_w\land T^K_w)\land T^K_w,\,
          T^K_w,T^K_w).
\]
It has two intervals and three required points. On ordered endpoints
it has a realization exactly when a strictly intermediate point exists.
This is precisely the same existential condition. In particular the
empty-family case is not an always-true interval formula.

If the outer endpoints are adjacent, there is no such intermediate
point, so every instance of \(\Gamma\) is false there. Every output
word has at least two intervals and thus requires at least one
interior point, so its disjunct is false there too. A singleton order
has no ordered endpoint pair and no strict midpoint; \(\Gamma\)
is false at its unique assignment and the property-\(I\) assertion
is vacuous under its guard. Global first or last points cause no
additional obstruction: the required midpoint need only lie between
the particular selected outer points. Density and endpointlessness
are not assumptions.

\paragraph{The one-interval example and finite termination.}
For the words \((A,B,C)\), \((D,E,F)\), the join definition
keeps the first gap \(B\), the second gap \(E\), and puts
\(H,C,D\) together at the intermediate point. Expanding the
two-interval basic past recursion gives exactly the body's displayed
formula, after associating its conjunctions and choosing fresh
bound names. In \(0<1<2\), the only possible midpoint between 0
and 2 is 1. All endpoint predicate requirements in the example
hold, and both open gaps are empty. Consequently the formula's
truth is exactly that of \(H\) at 1. For a concrete counterexample,
take \(H(w)=P_i(w)\), where \(P_i\) is a unary letter different
from \(A,B,C,D,E,F\), and interpret \(P_i\) by \(\{0,2\}\).
The composition is false while endpoint tests for \(H\) would
be true. This verifies the importance of the actual midpoint argument.

There are finitely many input conjunction steps. Each uses the finite
fusion construction of Lemma~4, and each final list is therefore finite.
Their finite pairwise product is explicitly enumerated. Every pair
uses a finite word concatenation and finite proper substitutions,
so the whole construction terminates. The joined point label uses
only already supplied elementary labels. Fresh skeleton letters
disappear on restoration, and the only extra fixed letter permitted
by the stated empty conventions is the source's \(P_0\).
This completes the positive-size source lemma and the explicitly
defined empty-family cases. Lemma~6 and its further pattern argument
have not been used.
\end{proof}
\end{annexitem}
